\section*{Chapter 5: Conclusion and Further Work}
\phantomsection \setcounter{section}{5} \setcounter{subsection}{0}
\addcontentsline{toc}{section}{\textit{Chapter 5: Conclusion and Further Work}}

\subsection{Conclusion}

This project examined a focused systems question: whether the control-plane invariants of MCP-style tool invocation are better enforced in the kernel than in a user-space gateway alone. 
The design, implementation, and evaluation of \texttt{linux-mcp} together support an affirmative answer.

The main outcome of the project is not simply a working prototype, but a more precise architectural argument.
Rather than treating AI-agent tool use purely as an application-layer integration problem, \texttt{linux-mcp} shows that a small but consequential part of the control plane benefits from operating-system-level support.

The design partitions control-plane state rather than reducing the trusted footprint of the gateway.
\texttt{mcpd} remains fully trusted: it populates the kernel's registries, computes semantic hashes, and derives the session-bound credentials used for identity verification.
What the kernel provides is a stronger enforcement domain for the state that \texttt{mcpd} has established: registry entries, approval tickets, and binding checks are protected by \texttt{CAP\_NET\_ADMIN}-enforced Netlink access control, survive daemon failure, and remain inspectable independently of the daemon's process lifetime.
Semantic parsing, schema validation, session logic, and tool execution remain in user space, where flexibility and evolvability are more important than privilege.

This partitioning is the central result of the project.
It shows that stronger control does not require relocating the entire agent stack into the kernel.
A narrow kernel role is sufficient to make the most critical state transitions---semantic identity, requester binding, approval validity, and invocation completion---explicit, durable, and independently inspectable.
In this sense, \texttt{linux-mcp} does not argue for a kernel-heavy AI operating system; it argues for careful, precisely scoped placement of a small set of control-plane invariants.

The evaluation further shows that this design choice is practically viable.
The kernel-assisted path blocks the tested classes of control-plane attack that continue to affect both the plain and hardened user-space baselines, while also preserving approval state across daemon failure in a way that neither baseline can match. 
These gains come at low cost: latency remains close to the user-space baselines, and throughput stays within the same practical range under the tested multi-agent workloads.

Taken together, these results suggest a broader conclusion. AI-agent systems do not necessarily require a fully redesigned operating-system substrate in order to obtain meaningful systems-level guarantees. 
A more grounded and deployable approach is to identify the invariants that matter most for safety and correctness, and to place only those in a stronger enforcement domain.
\texttt{linux-mcp} is therefore both a working prototype and a design claim: for MCP-style tool invocation, kernel-assisted control can improve integrity and resilience without sacrificing the extensibility that user space provides.

\subsection{Reflection}

The most important question clarified by this project was not simply "where should code run?", but "which invariants must remain trustworthy even when a gateway crashes, is misconfigured, or behaves incorrectly?"
That shift in perspective shaped the entire design.

From a technical standpoint, the project required me to work across layers that are often studied separately: implementing a Linux kernel module, using Generic Netlink as a structured kernel--user-space control interface, exporting live state through \texttt{sysfs}, and coordinating these pieces with a manifest-aware daemon and UDS tool services. 
This process deepened not only my implementation skills but also my understanding of interface ownership, failure handling, recovery logic, and the practical cost of crossing privilege boundaries.
More generally, it reinforced a central lesson of systems work: the harder problem is often not adding mechanisms, but deciding what each layer should---and should not---be responsible for.

A second lesson concerned research discipline. Early versions of the project could easily have been framed around broad claims such as "putting security in the kernel" or "making agents safer". 
In practice, such framing is too broad to be convincing. 
The project became much stronger once the problem was narrowed to a specific set of control-plane properties: semantic identity binding, approval validity, hash consistency, and durable observability. 
Once those invariants were explicit, the threat model, implementation choices, and evaluation criteria aligned much more naturally. 
This was an important lesson for me: a systems argument becomes credible only when the claimed benefit is specific enough to be tested directly.

The project also reinforced the value of restraint in privileged-system design.
A natural temptation is to move more logic into the kernel in order to claim stronger protection. 
In this case, that would have been the wrong direction.
Parsing manifests, interpreting rich schemas, and executing tools all belong in user-space, where change is frequent and debugging is tractable. 
Keeping the kernel role small was therefore not just an engineering convenience; it was part of the design argument itself. 
A privileged component that does too much becomes harder to reason about and, ultimately, harder to trust.

These technical choices also connect to broader ethical and professional issues.
Tool-using AI agents are not only language systems; they are action-taking systems. 
Once an agent can invoke external capabilities, failures in mediation may affect privacy, data integrity, or real-world operations. 
Approval correctness, auditability, and graceful failure handling are therefore not secondary details, but part of responsible system design. 
Throughout the project, I tried to keep the threat model precise, avoid overstating what the prototype guarantees, and distinguish clearly between the attack classes the system addresses and those that remain out of scope.

This project also made me more cautious about how such work should be presented.
It is easy to make strong claims around AI safety or operating-system support without equally strong evidence. \texttt{linux-mcp} improves a specific part of the control plane; it does not solve the full security problem of AI agents, and it does not address unsafe tool semantics, poor policy choices, or compromise outside the assumed threat model. Stating those limits clearly is part of good research practice.

Finally, the project highlighted a practical engineering principle that I expect to remain useful beyond this dissertation: better security is not always achieved by adding more machinery. 
In this case, the most effective step was to isolate a small set of authority-bearing state transitions and move them to a stronger enforcement domain. 
That is a modest conclusion, but a precise one, and it suggests that operating-system support for AI agents may be most valuable when it is targeted, limited in scope, and built around explicit invariants rather than ambitious but diffuse redesign.

\subsection{Further Work}

Several directions follow naturally from the current prototype.

First, the current kernel arbitration logic is intentionally minimal. 
It verifies tool identity, agent identity, semantic-hash consistency, and approval-ticket validity, but it does not implement a richer policy framework.
A natural next step is to support more expressive admission policies, such as per-tool permissions, role-scoped access control, rate limits, execution budgets, or policies conditioned on risk tags and request context. 
The challenge will be to extend policy expressiveness without turning the kernel into a semantic decision engine.

Second, the approval mechanism can be developed into a more complete governance path. 
The present design supports defer-and-decide control for sensitive operations, but real deployments may require approval delegation, bounded approval scope, revocation, expiry policies, and richer audit trails linking each decision to a concrete execution context. 
Strengthening this path would make \texttt{linux-mcp} more suitable for environments with accountability or compliance requirements.

Third, recovery behaviour can be improved further. The current prototype already preserves kernel-held registry and approval state across daemon failure, but reconciliation is deliberately conservative and session-level state is rebuilt after restart. 
A more mature design could provide more systematic restart handling, richer recovery metadata, and clearer support for interrupted or partially completed requests.
This becomes particularly important for long-running agent workflows, where failure and restart are routine rather than exceptional.

Fourth, the evaluation can be extended substantially. The current experiments
support the main claim of the project, but they remain limited in scale and
environment. Future work should include bare-metal evaluation, heavier concurrency, larger tool registries, and more heterogeneous tool workloads, including tools whose execution time is variable or network-dependent.
Comparison against additional hardened user-space designs would also help sharpen the argument by separating the benefits of kernel placement from those achievable through improved user-space engineering alone.

Finally, this work opens a broader research question: which other agent-oriented abstractions deserve operating-system support? 
This project focuses on admission control for tool invocation, but similar arguments may extend to execution budgets, provenance across tool chains, durable approval logs, or agent-scoped resource accounting. 
Exploring these directions would help clarify whether support for AI agents should be understood as entirely new operating-system substrates, or as conventional systems extended with a small number of carefully scoped, agent-aware mechanisms. 
In that sense, the contribution of \texttt{linux-mcp} is best read as a first step rather than a complete answer.

\vspace{\baselineskip}
\noindent\textbf{NOTE: The total length of the report up to this point should be between 30 and 50 pages.}
